\documentclass{ximera}

\title{Mean Value Theorem Activity}
\author{MATH 425: Calculus I}

\begin{document}
\begin{abstract}
   Working with the peers in your group, solve the following problems. Make sure to show and justify all your work. Make sure everyone in the group understands the solution and participates. Be prepared to report your answers to the whole class. 
\end{abstract}
\maketitle


\begin{exercise}
   Consider the following functions on the interval $[-2,4]$.  If the function satisfies the conditions of the Mean Value Theorem (MVT), find the $c$ as in the MVT.  If the function does not satisfy the conditions of the MVT, state why not.  
    \begin{enumerate}
       \item $f(x)=x^2-3x+4$\\
      %   {\color{blue} $f(4)=16-12+4=8$ and $f(-2)=4+6+4=14$\\
      %   $\frac{8-14}{4-(-2)}=-1$\\
      %   $f'(c)=2c-3$\\
      %   $2c-3=-1$ when $c=1$}
       \item $g(x)=\frac{1}{(x-4)^3}$\\
      %{\color{blue} $g(x)$ is not continous on $[-2,4]$ (infinite discontinuity at $x=4$). }
       \item $h(x)=|x-2|$\\
      %{\color{blue} $h(x)$ is not differentiable at $x=2$.}
       \item $D(x)=\ln(x+5)$\\
      %   {\color{blue} $f(4)=\ln(1)=0$ and $f(-2)=\ln(3)$\\
      %   $\frac{0-\ln(3)}{4-(-2)}=\frac{\ln(3)}{6}$\\
      %   $f'(c)=\frac{1}{c+5}$\\
      %   $\frac{1}{c+5}=\frac{\ln(3)}{6}$ cross multiply to get $6=\ln(3)(c+5)$\\
      %   $c\ln(3)=6-5\ln(3)$ or $c=\frac{6-5\ln(3)}{\ln(3)}$}
    \end{enumerate}

\end{exercise}

\begin{exercise}
   Ron’s toll pass recorded him entering the highway at mile 0 at 12:17 PM. He exited at mile 115 at 1:52 PM, and soon thereafter he was pulled over by the state police. “The speed limit on the highway is 65 miles per hour,” the trooper told Ron. “You exceeded that by more than five miles per hour this afternoon.” “No way!” responded Ron. “I glance at the speedometer frequently, and not once did it read over 65!” How did the trooper use the Mean Value Theorem to support her claim that Ron must have gone more than 70 miles per hour at some point? 

(Problem from Rogawski et al., \textit{Calculus: Early Transcendentals}, 4th Edition.)
\end{exercise}

\begin{exercise}
    Find a value of c that corresponds to the result of the MVT for the following functions on the given interval.  Sketch a graph of the tangent line to the function at the value you find, along with the secant line through the endpoints of the interval on the function.
   \begin{enumerate}
       \item $f(x)=x^5$, $[-1,1]$
       \item $g(x)=x\ln(x)$, $[1,2]$
   \end{enumerate}
\end{exercise}





\end{document}
